%% kapa: exam=1 round=36 year=2025 period=2 form=A subject=accounting subject_ko=회계학 seq=13 number=53
%% source: Q-Net 감정평가사 기출문제 (cst00309) artlSeq=5231525 file=raw/1008/5231525/2025년 제36회 감정평가사 1차시험 2교시 문제지 원본.pdf
%% conversion: kordoc 4.24.51 -> tools/gen1.py (auto); review=proofread-sonnet; flags=none
\begin{question}{53}
(주)감평은 20×1년 초 영업용 차량운반구(취득원가 ￦300,000, 내용연수 5년, 잔존가치 ￦50,000, 정액법 상각)를 취득하고 원가모형을 적용하였다. (주)감평은 20×2년 초 동 차량운반구를 매각하기로 결정하고 매각예정비유동자산으로 분류하였다. 동 차량운반구의 20×2년 초와 20×2년 말의 순공정가치와 사용가치가 다음과 같을 경우 (주)감평이 20×2년 말 재무상태표에 인식할 매각예정비유동자산의 장부금액은? (단, 동 차량운반구는 20×2년 초부터 20×2년 말 현재까지 매각예정 분류기준을 충족하고 있다.)
\par{\small\centering
\begin{tabular}{|c|c|c|}\hline
 & 20×2년 초 & 20×2년 말 \\ \hline
순공정가치 & ￦200,000 & ￦190,000 \\ \hline
사용가치 & 150,000 & 170,000 \\ \hline
\end{tabular}
\par}
\begin{choices}
  \item ￦150,000
  \item ￦170,000
  \item ￦180,000
  \item ￦190,000
  \item ￦200,000
\end{choices}
\end{question}
\answer{4}
